% Alignment's Diversity Cost Is Carried by a Small Part of the Shift
% arXiv submission source. Single file, no external .bib, no figures.
% Compile: pdflatex main && pdflatex main
\documentclass[11pt]{article}

\usepackage[margin=1in]{geometry}
\usepackage{amsmath,amssymb}
\usepackage{booktabs}
\usepackage[protrusion=true,expansion=false]{microtype}
\usepackage[utf8]{inputenc}
% NOTE ON FONTS: deliberately no \usepackage[T1]{fontenc} and no lmodern.
% T1 without cm-super/lmodern installed yields Type 3 bitmap fonts, which arXiv
% rejects. Default OT1 Computer Modern is Type 1 (scalable) everywhere. If you
% later add T1, add \usepackage{lmodern} in the same edit, never alone.
% Verify with: pdffonts main.pdf  -- every row must read "Type 1".
\usepackage[hidelinks]{hyperref}

% microtype pulls in textcomp, which routes \textbullet (the itemize marker) and
% \textsection through the TS1 encoding. With no cm-super/lmodern installed TS1
% has no Type 1 font, so those two glyphs come out as Type 3 bitmaps -- which
% arXiv rejects, and which pdffonts reports as a "Custom"-encoded Type 3 row.
% Route both back to Computer Modern math glyphs, which are Type 1.
% If you add a new symbol command and pdffonts grows a Type 3 row, it is the
% same cause: find the glyph with pdffonts, then map it here.
\AtBeginDocument{%
  \renewcommand{\textbullet}{\ensuremath{\bullet}}%
  \renewcommand{\textsection}{\ensuremath{\mathsection}}%
}

\newcommand{\rreg}{\rho_{\mathrm{reg}}}

\title{Alignment's Diversity Cost Is Carried by a Small Part of the Shift:\\
\large Stage-resolved measurement across four post-training pipelines}

%% ---------------------------------------------------------------
%% AUTHOR BLOCK
%% The postal address (Chandivali, Mumbai 400072) is deliberately NOT here.
%% arXiv collects it on the submission form, where it stays private; the PDF is
%% permanent and public, and a home street address in it cannot be withdrawn.
%% City and country are enough for a listing. Add the full address only if you
%% decide you want it public.
%% ---------------------------------------------------------------
\author{Amit Das\\
\small Mumbai, India\\
\small \texttt{amitdas@gmail.com}}
\date{\today}

\begin{document}
\maketitle

\begin{abstract}
Post-training alignment reduces the diversity of language model outputs, and the
loss is widely attributed to reinforcement learning on human preferences, with
\emph{typicality bias} --- a reweighting toward what the base model already found
likely --- offered as its mechanism. We measure both against four pipelines that
publish every intermediate checkpoint: OLMo-2-0425-1B, OLMo-2-1124-7B,
OLMo-2-1124-13B and Llama-3.1-T\"ulu-3-8B. Throughout, \emph{diversity} means
next-token entropy at the scored positions.

\paragraph{Composition and consequence come apart.} A typicality-aligned component
of each stage's log-probability shift accounts for a minority of that shift's
variance: a median of at most 2.2\% for instruction tuning and 7.1\% for the
preference stage under an unweighted fit allowing linear and piecewise-linear
terms, and at most about a third under a probability-weighted linear fit. An exact
counterfactual Shapley attribution of each stage's entropy change assigns that
same component \textbf{nearly all (87--139\%) of instruction tuning's entropy loss
and all (102--291\%) of the preference stage's}, under the pre-registered weighted
fit. Across six weighted specifications these become 77--157\% and 90--346\%;
under an unweighted fit they fall to $-$105--56\% and 35--178\%, so the magnitude
is a property of the weighting. Comparing like with like, a component accounting for
at most about a third of the shift carries nearly all of its entropy cost. The diversity cost is carried by a
small part of the shift.

\paragraph{By correlation, the larger of the two components we measured is a shared direction.}
Writing $\rreg$ for the
correlation between a stage's per-token shift and the mean shift on \emph{other}
prompts (leave-one-prompt-out), and $\psi$ for its correlation with reference
log-probability, $\rreg$ exceeds $\psi$ across the plausible set at every stage of
all four models, with intervals excluding zero and null controls at $+0.008$ and
$+0.004$. The preference-stage direction promotes dashes, em-dashes, ellipses and
brackets and demotes \emph{you} in all four. Neither component is most of the
shift: squared, the shared direction accounts for roughly 1--31\% of a stage's
shift and under 1\% for the 1B's instruction tuning, so the majority is explained
by neither axis. A second, distinct quantity --- how
far a stage's entropy falls on aligned-register text, which we call
\emph{register-conditional concentration} --- transfers across lineage: OLMo-1B's
instruction-tuning stage concentrates on text written by OLMo-2-1124-7B-Instruct
exactly as much as on its own family's ($-0.194$ both).

\paragraph{Where you measure changes the answer.} Scoring the same checkpoints on
the same prompts, T\"ulu-3-8B's preference stage accounts for 29.2\% [22.4, 35.0]
of total entropy loss at a chat-template position and 52.8\% [50.5, 55.4] over
generated continuations. Prompt-final scoring, closest to common practice, gives
confidence intervals 4.5--6$\times$ wider because it averages one position per
prompt rather than $\sim$190, and on the 1B model it falls 18--28 points below
every other convention. We have not found a published stage comparison that
reports this choice.

We also report two decoding-time interventions and a correction to our own earlier
work. Extrapolating along the base--aligned logit axis recovers diversity at a
third of the fluency cost of temperature at 135M (3.00$\times$, 95\% CI
[1.58, 5.14]) and does not survive to 1B (1.04$\times$, [0.73, 1.50], $n = 12$,
Wilcoxon $p = 0.97$). Reference subtraction, which an earlier draft reported as
catastrophic, is not: at matched diversity it sits at 0.98$\times$ the temperature
frontier, and the earlier figure compared it against $T = 1.0$ instead.
Plausible-set $\psi$ never exceeds 0.281 across four models, four stage deltas and
three conventions.
\end{abstract}

\section{Introduction}

That language models homogenise expression has moved from commentary into
measurement. Doshi and Hauser \cite{doshi2024} found that writers given LLM story
ideas produced work rated more creative individually while the corpus became more
similar in aggregate. Sourati, Ziabari and Dehghani \cite{sourati2026} synthesised
evidence across four disciplines that models reinforce dominant styles.

The technical literature localises the loss to the alignment step. Kirk et al.\
\cite{kirk2024} established a generalisation--diversity trade-off under RLHF, and
the phrase ``RLHF reduces diversity'' is standard. \emph{Typicality bias}
\cite{zhang2025} supplies a mechanism: human preference data favours the
conventional, and KL-regularised optimisation against it amplifies majority-style
responses.

Both are usually evaluated by comparing a base model with a released instruct
model, collapsing instruction tuning, preference optimisation and any later stages
into one delta. Pipelines publishing intermediate checkpoints allow the stages to
be separated. We separate them on four, and ask two questions of each stage: what
is its shift made of, and which part of it removes entropy. The answers are
different components.

\subsection{Contributions}

\begin{enumerate}
\item \textbf{Composition is not consequence.} In every model and at both stages, a
typicality-aligned component accounting for a minority of the shift's variance
carries nearly all of its entropy loss.
\item \textbf{A shared-shift geometry, measured and replicated.} $\rreg > \psi$
across the plausible set at every stage of all four models, under two continuation
sources, with null controls near zero and the token direction stated explicitly.
In the probability-weighted region the ordering reverses (Section~\ref{sec:head}).
\item \textbf{Cross-lineage transfer of register-conditional concentration},
ruling out checkpoint ancestry though not shared training data.
\item \textbf{The measurement convention determines the stage attribution.} Four
conventions on four models; the same checkpoints yield preference-stage shares 24
points apart, and the effect is one of statistical precision as well as bias.
\item \textbf{Two decoding-time negative results}, one scale-dependent, plus
corrections to a matched-diversity error and a withdrawn threshold in our own
earlier drafts.
\end{enumerate}

Every claim below is scoped to the models, conventions and regions it was measured
on. None is stated as general across post-training pipelines.

\section{Measurement}

\subsection{Models and protocol}

\begin{center}
\begin{tabular}{llll}
\toprule
pipeline & base & stages & params \\
\midrule
OLMo-2-0425 & OLMo-2 & base $\to$ SFT $\to$ DPO $\to$ Instruct (RLVR) & 1B \\
OLMo-2-1124 & OLMo-2 & base $\to$ SFT $\to$ DPO $\to$ Instruct (RLVR) & 7B \\
OLMo-2-1124 & OLMo-2 & base $\to$ SFT $\to$ DPO $\to$ Instruct (RLVR) & 13B \\
Llama-3.1-T\"ulu-3 & Llama 3.1 & base $\to$ SFT $\to$ DPO $\to$ Instruct (RLVR) & 8B \\
\bottomrule
\end{tabular}
\end{center}

The 7B and 13B share a release, so comparisons between them vary scale alone. The
1B is a different release (April 2025 vs November 2024), so comparisons involving
it hold the recipe near-constant rather than constant.

\textbf{Diversity, in this paper, is next-token entropy at the scored positions.}
The generation experiments of Section~\ref{sec:decode} use mean pairwise Jaccard
instead, and say so.

36 prompts across six task families. Each checkpoint is loaded alone, in bf16 with
a per-parameter assertion after load, and its next-token log-probability vectors
cached before the next is loaded; all delta arithmetic happens afterward on saved
arrays. Confidence intervals are bootstrap over prompts, 20{,}000 resamples.

The stage, register and decomposition analyses ran on Apple MPS; the 1B generation
sweep of Section~\ref{sec:lambda} and all 135M runs ran on CPU. The T\"ulu pipeline
was reproduced in full on CPU under the prompt-final convention and agrees on every
share to within 0.3 percentage points.

Tokeniser identity was verified at the token-ID level for every pipeline, a
requirement for logit arithmetic rather than a formality. OLMo needed no
restriction. T\"ulu's SFT, DPO and final heads are 128{,}264 wide against
Llama-3.1-8B's 128{,}256, so T\"ulu comparisons use the shared 128{,}256 IDs,
renormalised; excluded mass reaches $7\times10^{-6}$ at prompt-final positions,
$1.5\times10^{-4}$ over base-written continuations and $7.8\times10^{-4}$ over
final-written ones.

\subsection{Four conventions}
\label{sec:fourconv}

Entropy and the per-token shift must be evaluated at some set of positions, and the
literature does not state which.

\begin{itemize}
\item \textbf{Raw prompt-final.} One position per prompt: the last token of the raw
prompt. Closest to common practice and what our own earlier drafts used.
\item \textbf{Chat-templated.} The prompt wrapped in the aligned model's chat
template, scored at the assistant-start position, with identical token IDs fed to
all four stages.
\item \textbf{Continuation, base-written (primary).} The base model generates
$6\times32$ tokens per prompt; all four stages are teacher-forced through that text
so contexts are paired. $\sim$190 positions per prompt.
\item \textbf{Continuation, final-model-written.} The same, with the final model as
writer.
\end{itemize}

Base-written continuations were pre-registered as the source for the register test
of Section~\ref{sec:geom} before any of that data existed. They were adopted as
primary for the stage shares later, in a dated note appended after the 1B, 7B and
8B continuation shares had been seen but before any 13B data was collected. Both
pre-registrations are released with their hashes at \url{https://github.com/amitd01/axisalignment-paper}, and the appended
note is part of the record.

\subsection{The convention changes the answer}
\label{sec:conv}

Preference-stage share of total entropy loss, by convention:

\begin{center}
\small
\begin{tabular}{lcccc}
\toprule
model & prompt-final & chat-template & continuation (base) & continuation (final) \\
\midrule
OLMo-2-1B & $\mathbf{-0.6\%}$ [$-15.4$, 9.1] & 25.5\% [18.0, 35.2] & 27.6\% [24.1, 31.2] & 17.4\% [15.1, 20.1] \\
OLMo-2-7B & 25.7\% [19.9, 31.9] & 26.1\% [21.9, 30.5] & 33.3\% [30.9, 35.8] & 23.9\% [22.0, 25.9] \\
T\"ulu-3-8B & 40.8\% [33.0, 50.0] & \textbf{29.2\%} [22.4, 35.0] & \textbf{52.8\%} [50.5, 55.4] & 44.3\% [41.8, 47.0] \\
OLMo-2-13B & 31.4\% [22.3, 38.7] & 26.8\% [22.0, 31.6] & 49.0\% [46.0, 52.2] & 39.7\% [37.1, 42.4] \\
\bottomrule
\end{tabular}
\end{center}

T\"ulu is the clearest case: 29.2\% and 52.8\% are the same checkpoints and the
same prompts, 24 points apart, and neither position is obviously the wrong one to
score.

\paragraph{Prompt-final breaks at small scale.} On the 1B model it sits 18--28
points below every other convention, reporting that preference optimisation does
nothing measurable --- a conclusion no other convention supports. At 7B, 8B and 13B
the raw value falls inside or near the range of the others. We observe this at
$n = 1$ small model and do not generalise it.

\paragraph{The spread is not monotone in scale:} 28.2 points at 1B, 9.4 at 7B, 23.6
at 8B and 22.2 at 13B. The 7B and 13B share a release and differ by more than a
factor of two, so sensitivity to position is not a fixed property of a pipeline
either.

\paragraph{The effect is about precision as much as bias.} Prompt-final averages one
position per prompt; continuation scoring averages $\sim$190. Total-loss confidence
intervals are correspondingly 4.5--6$\times$ wider by model. Comparing OLMo-7B
against T\"ulu-8B: under prompt-final the totals read 1.123 and 0.963, a paired
difference of 0.160 [$-0.016$, 0.325] --- indistinguishable; under continuation
scoring they read 0.765 and 0.756, difference 0.009 [$-0.038$, 0.060] --- identical
to within $\pm0.05$ nats. The convention moves both the point estimates and how
much can be concluded from them.

\paragraph{Whose text matters too, through the denominator.} Scoring the final
model's own text rather than base-written text lowers the preference share by 8--10
points in every model. This is not the preference stage behaving differently ---
its absolute drop barely moves ($-0.01$ to $-0.04$ nats) --- but instruction tuning
doing more ($+0.07$ to $+0.19$ nats), enlarging the total. The obvious explanation,
that the base model is out of distribution on aligned text, is ruled out: base
entropy is not higher there but lower at 7B, 8B and 13B ($-0.21$, $-0.27$, $-0.14$)
and unchanged at 1B ($-0.055$ [$-0.112$, $+0.008$]).

\paragraph{Reporting standard.} State the position convention. Report every stage
share beside its absolute drop in nats and the pipeline total, because the three
move independently --- Section~\ref{sec:stages} gives a case where a share rises
while its absolute cost falls and the total halves.

\section{Stage attribution}
\label{sec:stages}

Under the primary convention:

\begin{center}
\begin{tabular}{lcccc}
\toprule
model & SFT & preference & RLVR & total loss \\
\midrule
OLMo-2-1B & \textbf{66.1\%} [62.2, 69.7] & 27.6\% [24.1, 31.2] & 6.3\% [5.0, 7.7] & 0.426 \\
OLMo-2-7B & \textbf{60.2\%} [57.5, 62.8] & 33.3\% [30.9, 35.8] & 6.5\% [6.0, 7.0] & 0.765 \\
T\"ulu-3-8B & 40.7\% [37.9, 43.2] & \textbf{52.8\%} [50.5, 55.4] & 6.5\% [5.8, 7.3] & 0.756 \\
OLMo-2-13B & 40.5\% [36.7, 44.1] & \textbf{49.0\%} [46.0, 52.2] & 10.5\% [9.8, 11.4] & 0.413 \\
\bottomrule
\end{tabular}
\end{center}

Which stage is larger varies by model, and within the OLMo-1124 release it flips
between the 7B and the 13B. The universal form of ``RLHF reduces diversity'' is not
supported, and neither is its negation.

\paragraph{A share is not a cost.} Between the 7B and 13B --- same release, scale
the only variable --- the preference share rises from 33.3\% to 49.0\%, a paired
difference of $+15.6$ points [$+12.7$, $+18.6$]. Over the same step its
\emph{absolute} entropy cost falls, $0.255 \to 0.202$ nats ($-0.053$
[$-0.078$, $-0.028$]), and the pipeline total nearly halves, $0.765 \to 0.413$,
because instruction tuning's cost collapses from 0.460 to 0.167. Under the
chat-template convention nothing changes at all ($+0.7$ [$-3.6$, $+4.7$]). ``The
preference share rises with scale'' is true here and means close to the opposite of
what it sounds like.

Total collapse is non-monotone within OLMo ($1\mathrm{B} < 7\mathrm{B} >
13\mathrm{B}$ under both prompt-final and continuation scoring). The 13B loses the
least of any model under prompt-final (0.469 against 0.602 for the next lowest),
but under continuation scoring it is indistinguishable from the 1B ($-0.013$
[$-0.053$, $+0.027$]), and under the chat template it loses more. There is no scale
trend in total collapse across these four.

RLVR is small everywhere: 3.6--10.5\% across all models and conventions, though the
7B and 13B differ within the primary convention (6.5\% [6.0, 7.0] against 10.5\%
[9.8, 11.4]).

\section{What the shift is made of}
\label{sec:geom}

\subsection{Two geometries}

For each stage delta we compute the per-token shift
$\delta(v) = \log p_{\mathrm{target}}(v \mid c) - \log p_{\mathrm{ref}}(v \mid c)$
and correlate it with two things:

\begin{itemize}
\item $\psi$ --- the reference model's own log-probability,
$\log p_{\mathrm{ref}}(v \mid c)$. This is the quantity the typicality account
makes a claim about.
\item $\rreg$ --- the mean shift at all positions of the \emph{other} 35 prompts
(leave-one-prompt-out, so no in-sample leakage). High $\rreg$ means the stage
applies the same shift regardless of context.
\end{itemize}

We also report the partial correlation of $\delta$ with the shared direction
controlling for $\log p_{\mathrm{ref}}$, which asks whether the shared direction is
typicality under another name.

\subsection{The shared shift exceeds typicality across the plausible set}
\label{sec:exceeds}

Pre-registered test, base $\to$ SFT, plausible region
($p_{\mathrm{ref}} > 10^{-6}$), base-written continuations:

\begin{center}
\begin{tabular}{lcccc}
\toprule
model & $\psi$ & $\rreg$ & partial & $\rreg - \psi$ \\
\midrule
OLMo-2-1B & $-0.017$ & $\mathbf{+0.092}$ [$+0.083$, $+0.101$] & $+0.089$ & $+0.109$ [$+0.098$, $+0.120$] \\
OLMo-2-7B & $+0.082$ & $\mathbf{+0.320}$ [$+0.310$, $+0.331$] & $+0.324$ & $+0.238$ [$+0.222$, $+0.253$] \\
T\"ulu-3-8B & $+0.069$ & $\mathbf{+0.263}$ [$+0.248$, $+0.276$] & $+0.267$ & $+0.194$ [$+0.178$, $+0.208$] \\
OLMo-2-13B & $+0.118$ & $\mathbf{+0.445}$ [$+0.431$, $+0.457$] & $+0.447$ & $+0.326$ [$+0.307$, $+0.345$] \\
\bottomrule
\end{tabular}
\end{center}

All four pass the pre-registered criterion: $\rreg > 0$, $\rreg > \psi$,
partial $> 0$, all intervals excluding zero, under both continuation sources (the
smallest margin under final-written text is 0.097 [0.085, 0.111], at the 13B
preference delta). \textbf{Null controls} --- identical weights evaluated on two
numerical backends, run on OLMo-1B and T\"ulu-8B --- give $\rreg = +0.008$
[$+0.003$, $+0.014$] and $+0.004$ [$-0.000$, $+0.009$]. The partial correlation is
within 0.003 of, or larger than, the raw one in every case, so the shared direction
is not typicality in disguise.

The preference delta shows the same pattern more strongly at the 1B, 7B and T\"ulu
($\rreg = +0.558$, $+0.489$, $+0.486$) and less so at the 13B, where it is $+0.386$
--- smaller than that model's base $\to$ SFT value, with $\rreg - \psi$ shrinking
from 0.326 to 0.130.

\subsection{The direction, in tokens}

The preference-stage shared shift promotes dashes, em-dashes, ellipses and
brackets, and demotes \emph{you}, in all four models and under both continuation
sources. \emph{I} is among the fifteen most-demoted tokens in the 1B, 7B and T\"ulu
under base-written text, but not in the 13B under either source, nor in T\"ulu
under final-written text. In the 1B and 7B the shift also demotes newlines and
end-of-sequence --- which instruction tuning had \emph{promoted} in those same
models --- while T\"ulu and the 13B promote paragraph-break tokens instead. Hedges
(\emph{effectively}, \emph{essentially}, \emph{clearly}) are top-ranked only in the
1B.

Token lists are per delta and per model; the stable core across all four is
punctuation up, second person down.

\subsection{A related quantity transfers across lineage}
\label{sec:lineage}

This subsection concerns a different quantity from the preceding one.
\textbf{Register-conditional concentration} is how much further a stage's entropy
falls when scoring aligned-register text rather than base-written text;
\textbf{shared shift} ($\rreg$) is the context-invariance of the log-probability
shift. The lineage test below measures the former, not the latter.

Register-conditional concentration could be nothing more than a model being sharper
on text written by its own descendants. OLMo-1B and OLMo-2-1124-7B share a
tokeniser, which permits a direct test: score 7B-Instruct-written text through the
1B stages, where no 1B checkpoint is an ancestor of the writer.

The concentration is identical: $\mathbf{X[\mathrm{SFT}] = -0.194}$ [$-0.224$,
$-0.163$] against same-lineage $-0.194$ [$-0.230$, $-0.158$]. Later stages add
nothing beyond SFT.

This rules out checkpoint ancestry as the explanation for register-conditional
concentration. It says nothing directly about whether the shared direction
transfers; we did not test that. It does not rule out shared training data --- both
pipelines use T\"ulu-3-style instruction mixtures --- and that remains the live
alternative explanation.

\subsection{At the head of the distribution it looks different}
\label{sec:head}

The picture above is of the broad plausible set, with a median of roughly
900--3{,}800 tokens per position. In the probability-weighted region the ordering
reverses: $\psi_w$ exceeds $\rho_w$ for base $\to$ SFT in all four models, and for
the preference delta at 7B, 8B and 13B (the 13B reads $\psi_w$ 0.433 against
$\rho_w$ 0.096; OLMo-7B's base $\to$ SFT gives $\rho_w - \psi_w = -0.203$).
OLMo-1B's preference delta is the exception, with $\rho_w$ 0.435 $>$ $\psi_w$
0.163.

The probability-weighted statistic is heavily concentrated: the tokens covering
50\% of base mass number a \emph{median of two}. So ``the head looks like
reweighting'' is a statement about the top one or two candidates per position, not
about half the distribution. Both readings hold at their own scope.

\paragraph{$\psi$ in absolute terms.} Across four models and four stage deltas,
plausible-set $\psi$ never exceeds \textbf{0.281} under prompt-final and the two
continuation conventions, and never exceeds 0.256 under the primary convention.
Probability-weighted, it never exceeds 0.566 under those three; including
chat-template positions it reaches 0.663 (13B, DPO $\to$ final), whose
plausible-set maximum is 0.237.

\section{What costs the diversity}
\label{sec:cost}

Composition and consequence are different questions, and this section separates
them. We attribute each stage's entropy change using counterfactual
distributions --- apply only part of the stage's log-probability shift,
renormalise, recompute entropy --- with Shapley values over the parts, which sum
exactly to the measured $\Delta H$. Shapley efficiency does not remove interactions;
it allocates them among the parts. This is an attribution under a stated
decomposition, not causal identification of what training did.

Two terms are used narrowly. The \emph{typicality-aligned component} is the fitted
term $\beta \cdot \log p_{\mathrm{ref}}$, a reference-probability-aligned part of the
shift; reading it as the typicality bias of the source literature is an
interpretation (Sec.~\ref{sec:limits}, limitation 4). The \emph{shared-shift
component} is $\gamma \cdot g$, where $g$ is the leave-one-prompt-out mean shift.
Both are fitted at every position, so the bootstrap needs no refit; $g$ itself is
computed once on the full sample and held fixed across resamples.

\subsection{A component explaining a minority of the shift carries nearly all of the loss}
\label{sec:minority}

\paragraph{The pre-registered result.} Under the weighted linear fit specified in
advance, the typicality-aligned component carries \textbf{nearly all of instruction
tuning's entropy loss} (87.0\% [84.9, 88.9] at the 13B, 94.5\% [92.7, 96.4] at the
7B, 96.6\% [94.5, 98.8] at T\"ulu, 139.0\% [131.7, 147.5] at the 1B) \textbf{and
all of the preference stage's} (102.2\% to 290.5\%). Only the preference stage and
the 1B's instruction tuning reach 100\%; the other three instruction-tuning
intervals sit below it. Shares above 100\% with opposite-signed partners mean the
components interact; magnitudes are not separable there, but the sign and the
ordering are.

\emph{A worked example.} The 1B's preference stage loses 0.118 nats per position.
Averaged over the orders in which the parts can be added, the typicality-aligned
part removes 0.342 nats (290.5\%); the shared-shift part restores 0.278
($-$236.2\%); the remainder removes 0.054 (45.8\%) and the off-set tokens 0.000
($-$0.2\%). The four sum to the measured loss. A share above 100\% means the
component removes more entropy than the stage does in total, because another part
puts some back.

\paragraph{How much of the shift that component is --- exploratory.} The variance
analysis below was run after the other results, to ask how large the costly
component is. It was not pre-registered and is reported as exploratory. Median
$R^2$ of each stage's shift over the plausible set explained by typicality (linear
and piecewise-linear in reference log-probability), base-written continuations,
400 sampled positions per model:

\begin{center}
\begin{tabular}{lcc}
\toprule
model & base $\to$ SFT & SFT $\to$ DPO \\
\midrule
OLMo-2-1B & 0.7\% & 1.0\% \\
OLMo-2-7B & 1.5\% & 3.6\% \\
T\"ulu-3-8B & 1.4\% & 6.2\% \\
OLMo-2-13B & 2.2\% & 7.1\% \\
\bottomrule
\end{tabular}
\end{center}

Recomputed over all $\sim$6{,}800 positions per model with bootstrap intervals,
these hold: 0.73\% [0.68, 0.76], 1.57\% [1.37, 1.78], 1.45\% [1.29, 1.64] and
2.20\% [1.80, 2.58] for instruction tuning across the 1B, 7B, T\"ulu and 13B, and
1.1--7.1\% for the preference stage with an upper bound of 7.9\%. Under a
probability-weighted \emph{linear} fit --- the pre-registered one, with no
piecewise-linear term --- the same quantity is larger: up to 19.2\% [16.9, 21.4]
for instruction tuning and 30.4\% [27.96, 33.42] for the preference stage. An
exploratory variant merging linear and non-linear typicality assigns the component
53--122\% and 95--204\% of the two stages' losses.

\textbf{Comparing like with like --- the weighted linear fit on both sides --- a
component accounting for at most about a third of a stage's shift carries nearly all
of its entropy loss, in every model and at both stages.} Under an unweighted fit
admitting a piecewise-linear term, the same component accounts for at most 7.1\%.
That is the paper's central measurement, and the weighted pairing is the
conservative one. The weighting is not arbitrary: entropy is $-\sum p \log p$, so
each token's contribution to it is already weighted by its probability, and a fit
weighted the same way is judged on the tokens that move entropy. The unweighted fit
treats a token at $p = 10^{-6}$ like one at $p = 0.3$. Note also that the two halves
of the sentence have different standing: the entropy shares are pre-registered, the
variance shares exploratory.

\paragraph{Robustness --- exploratory.} \texttt{code/robustness.py} repeats the
decomposition over twelve specifications: weighted or unweighted fit; a plausible
set of $p_{\mathrm{ref}} > 10^{-5}$, $10^{-6}$ or $10^{-7}$; and a shared direction
estimated from all other prompts or only from prompts in other task families. It
reports the variance share and the entropy share under the same fit, with a
bootstrap over prompts stratified by task family. These runs score every fourth
position; at the pre-registered specification they reproduce the full-position
shares to within 2.1 points at the 7B, T\"ulu and 13B, and run 8.6 and 31.2 points
high at the 1B, whose loss is smallest (\texttt{logs/robustness\_s4.txt}).

\begin{center}
\small
\begin{tabular}{llrrr}
\toprule
stage & model & weighted fit & unweighted fit & weighted R$^2$ \\
\midrule
instruction tuning & 1B & 127.9 to 156.8\% & $-$104.8 to 35.2\% & 14.4 to 15.4\% \\
 & 7B & 83.8 to 100.8\% & 9.2 to 45.0\% & 24.7 to 25.2\% \\
 & T\"ulu-8B & 86.0 to 102.5\% & 5.1 to 52.1\% & 20.5 to 20.9\% \\
 & 13B & 77.3 to 90.8\% & 43.3 to 56.1\% & 20.5 to 20.8\% \\
preference & 1B & 266.8 to 345.5\% & 111.1 to 178.4\% & 15.7 to 16.4\% \\
 & 7B & 115.3 to 135.3\% & 35.1 to 80.6\% & 27.6 to 27.8\% \\
 & T\"ulu-8B & 90.8 to 111.4\% & 51.1 to 73.3\% & 33.8 to 34.2\% \\
 & 13B & 89.6 to 107.0\% & 65.8 to 82.0\% & 27.4 to 28.2\% \\
\bottomrule
\end{tabular}

\emph{Typicality-aligned share of each stage's entropy loss, minimum to maximum over
the six specifications of each fit; last column, median per-position variance share
under the weighted fit.}
\end{center}

Under the weighted fit the result survives every specification: the component
carries 77.3--156.8\% of instruction tuning's loss and 89.6--345.5\% of the
preference stage's while accounting for 14.4--34.2\% of the shift's variance.
``All'' of the preference stage's loss becomes ``nearly all'' at the loosest
threshold, where the 13B and T\"ulu move 16.5\% and 18.1\% of the loss to tokens
outside the plausible set. Holding the shared direction out by task family changes
no share by more than 3.4 points, so it is not a within-family formatting habit.
Within each of the six families the pre-registered share is 81.7--167.4\% and
99.0--495.7\%, and leaving one family out moves it by at most 4.4 points ---
except at the 1B preference stage, whose components nearly cancel. Under the
unweighted fit only the contrast survives: the entropy share exceeds the variance
share in 92 of 96 cells, but the share itself falls to $-$104.8--56.1\% and
35.1--178.4\%, and at the 1B's instruction tuning its sign depends on the threshold
(the four exceptions). \textbf{``Nearly all'' is a property of the
probability-weighted fit.} The final stage, reported without prediction, matches the
others at the 7B, T\"ulu and 13B (typicality 89.8--91.2\%) but not at the 1B, where
the shared direction carries 80.3\% and typicality 9.0\% [$-$15.2, 25.2].

\paragraph{The remainder is not inert.} For the preference stage, the component
orthogonal to both axes falls with scale within the OLMo-1124 release: 25.7\%
[22.1, 29.3] of the entropy loss at the 7B against 7.4\% [5.2, 9.7] at the 13B,
with T\"ulu-8B --- a different family --- at 14.7\% [12.2, 17.5]. The 1B's 45.8\%
[37.0, 57.0] is not comparable to those: its decomposition nearly cancels there
(typicality 290.5\% against a shared share of $-236.2\%$), so only the positive
sign is interpretable.

At instruction tuning the 1B is the interesting case, and it runs the other way:
the remainder \emph{raises} entropy, $-30.6\%$ [$-37.7$, $-24.1$], partly
offsetting a typicality component of 139.0\%. That decomposition does not
near-cancel --- the shared share is 2.1\% --- so unlike the 1B preference stage
its magnitude can be read. At the 7B, T\"ulu and 13B the remainder carries
1--6\%.

Attributing nearly all of the loss to the typicality-aligned component is
therefore a statement about that component's share, not a demonstration that
nothing else moves entropy. We do not characterise what the remainder is.

\subsection{Whether the two stages differ depends on the fit}
\label{sec:fitdep}

\paragraph{The preference stage.} Its entropy loss comes mostly from the body of the
distribution: the head (reference-top tokens covering 50\% of mass, median 2
tokens) carries 15--23\%, the body 74--86\%. The decomposition's pre-registration
predicted the opposite --- a head share above 50\% --- so its primary hypothesis is
\textbf{not supported} in any model; its secondary condition, typicality above
shared shift, holds in all four. Decomposing by component, the loss is
carried by the typicality-aligned part while the shared-shift part pushes entropy
back \emph{up} --- by 16\%, 28\% and 55\% at the 13B, T\"ulu-8B and OLMo-7B on the
weighted fit (18\%, 28\%, 47\% unweighted). The weighted typicality-minus-shared
contrast runs $+118$ to $+527$ points with intervals clear of zero in all four
models.

Sampling 400 positions each \textbf{at the 1B, 7B and T\"ulu-8B} (not the 13B), the
typicality-aligned component is a positive sharpening at 85--97\% of positions,
with weighted median $\beta \approx +0.09$ to $+0.12$ --- a local temperature of
about 0.9. The 1B is the exception \textbf{for the preference stage}:
$\beta \approx 0$ there, with components nearly cancelling, so only their signs are
interpretable. Its instruction-tuning stage behaves like the others,
$\beta \approx +0.10$. $\beta$ was not measured at the 13B, where only the Shapley
shares are available.

\paragraph{Instruction tuning, and a fit-dependence we initially mistook for a
mechanism.} An earlier draft reported that instruction tuning's costly component was
not a temperature while the preference stage's was. That contrast is an artifact of
comparing the two stages under different fits.

Under the \textbf{weighted} fit --- the pre-registered one --- the two stages look
alike. Instruction tuning's sharpening coefficient is $\beta \approx +0.08$ to
$+0.115$, against the preference stage's $+0.09$ to $+0.12$, and linear typicality
alone carries 87--139\% of instruction tuning's loss.

Under an \textbf{unweighted} fit they diverge. A single linear coefficient then
captures between $-29\%$ and 51\% of instruction tuning's loss, and a
piecewise-linear function of reference log-probability is needed: it captures
200\%, 54\%, 76\% and 23\% across the four models on the pre-registered statistic,
or 122\% / 56\% / 67\% / 53\% merging linear and non-linear typicality (they are
collinear, so the split between them is arbitrary). The 7B and 13B interval lower
bounds are 49\% and 41\%.

So: \textbf{under an unweighted fit instruction tuning's loss needs a non-linear
term and the preference stage's does not; under a weighted fit both look like a
local sharpening of similar size.} We report both and claim no mechanism difference
between the stages.

A further earlier claim --- that no decoding-time operator built on the reference
axis could reach instruction tuning's loss --- is \textbf{not supported}. On the
pre-registered non-linear statistic it is too strong at the 1B and at T\"ulu,
unresolved at the 7B, and holds only at the 13B. A non-linear, truncation-like
reference operator is not ruled out.

\section{Two decoding-time interventions (exploratory)}
\label{sec:decode}

The evidence in this section is thin by design: five prompts at 135M, twelve at
1B, and a single fluency judge (GPT-2, 124M). It is reported as counterexamples to
earlier predictions, not as validated methods for recovering diversity.

\subsection{Extrapolating along the base--aligned axis}
\label{sec:lambda}

DeRa \cite{liu2024dera} proves that aligned models at varying KL strength are
geometric mixtures differing only in mixing weight:
\[
z = z_b + \lambda (z_a - z_b)
\]
$\lambda = 1$ recovers the aligned model, $\lambda = 0$ the base model,
$\lambda < 0$ extrapolates past the base model away from alignment. Published
evaluations cover $\lambda \in [0, 2]$.

Diversity in this section is mean pairwise Jaccard, and fluency is median GPT-2
perplexity --- an independent judge sharing neither family nor tokeniser with
either model in the mixture.

At \textbf{135M} (SmolLM2 base/instruct), the per-prompt paired comparison at
matched diversity gives $\lambda = -1.0$ beating temperature by
\textbf{3.00$\times$, 95\% CI [1.58, 5.14]}, winning on 4/5 prompts. Re-reading the
published scored numbers at matched diversity gives $\lambda < 0$ ratios of
0.17--0.37 against the temperature frontier.

At \textbf{1B} it does not survive. With the full design of 12 prompts, the paired
comparison gives \textbf{1.04$\times$, 95\% CI [0.73, 1.50]}, $\lambda$ winning on
5/12, Wilcoxon $p = 0.970$ against an attainable floor of 0.0005. The estimate
declines monotonically from $n = 7$ onward: 1.16$\times$, 1.08$\times$,
1.04$\times$. The interval excludes the 135M effect size, whose lower bound is
1.58, but does not exclude a modest effect up to 1.5$\times$.

A direction control --- the delta shuffled across the vocabulary at each step,
preserving its norm and the full marginal distribution of its values --- is
decisive at 135M (401.3 against temperature's 403.1 at matched diversity, while the
true delta gives 165.4) and \textbf{uninformative at 1B}, where both arms have
degenerated into the thousands by the time their diversity matches.

\subsection{Reference subtraction, and a correction}
\label{sec:refsub}

Writing $p_\gamma \propto p_{\mathrm{align}} \cdot p_{\mathrm{ref}}^{-\gamma}$
attenuates the reference's contribution --- the repair the typicality account
recommends, and contrastive decoding \cite{li2023} with the base model as amateur.

An earlier draft of this work reported that $\gamma = 0.5$ raised judge perplexity
to 19{,}956 and treated this as catastrophic failure. \textbf{That comparison was
against $T = 1.0$, not at matched diversity.} Re-read at matched diversity,
$\gamma = 0.5$ sits at \textbf{0.98$\times$ the temperature frontier} at its own
Jaccard. It overshot along the frontier; it did no off-frontier harm. The original
matched-diversity table stopped short of that diversity level, so the condition was
never properly compared.

A pre-registered sweep $\gamma \in \{0.05, 0.1, 0.2, 0.3, 0.5\}$ on the 135M pair
finds nothing degenerate at small $\gamma$. $\gamma = 0.05$ is classified
temperature-like (1.21$\times$, CI [0.79, 1.33]) but fragile --- only 52\% of
bootstrap draws fall inside the temperature curve's interpolation range;
$\gamma = 0.1$ misses the band by 0.03. With five prompts the sweep neither
falsifies nor establishes the hypothesis.

The measured $\beta$ on that pair is $+0.081$ at continuation positions. That is the
\emph{whole} base $\to$ Instruct pipeline, since SmolLM2 publishes no intermediate
checkpoints, so it is comparable in size to the preference-stage $\beta$ above but
is not the same quantity.

\subsection{A partial explanation}

For the preference stage at 7--8B, the decomposition suggests why an operator on the
reference axis would track the temperature frontier rather than beat it: the
entropy-relevant part of that stage's shift is itself close to a temperature.

This cannot be the explanation for the 1B null. There the preference stage's
$\beta \approx 0$, and $\lambda$ acts on the whole base $\to$ Instruct delta, which
at the 1B is 66\% instruction tuning --- whose costly component looks non-linear
under an unweighted fit and temperature-like under a weighted one. We offer the
account as a hypothesis for 7--13B and note that the 1B result is unexplained.

\section{What we got wrong}
\label{sec:failures}

Four earlier conclusions in this work were overturned, by four different causes. We
list them with their causes because a vague account of one's own errors is worth
less than a precise one.

\begin{center}
% p-column on the right: the "what overturned it" cells vary a lot in length and
% an ll table overflows the text block as soon as one of them grows.
\begin{tabular}{l p{0.34\textwidth}}
\toprule
earlier claim & what overturned it \\
\midrule
``Instruction tuning carries $\sim$99\% of the entropy loss'' & position convention \\
A 1B $\to$ 7B jump in the preference share & position convention \\
``Reference subtraction fails catastrophically'' & comparison baseline ($T = 1.0$ instead of matched diversity) \\
``No reference-axis operator can reach SFT's loss'' & operationalisation (linear-only fit) \\
\bottomrule
\end{tabular}
\end{center}

Two further claims fell to ordinary causes: ``instruction tuning does all of it'' as
a \emph{general} claim was refuted by measuring a second pipeline, and a reported
negative sign on the 1B preference-delta $\psi$ did not survive a change of
numerical backend.

Two analytical predictions also failed, in opposite directions.

\paragraph{$\lambda < 0$ was predicted catastrophic.} An analytical model predicted
$\lambda < 0$ would be harmful at any scale, because $\lambda$ rescales the whole
reward and reversing it reverses the quality term with the typicality term. At 135M,
generation showed the opposite. The model treated the reference as an abstract
distribution over items with independently drawn quality; real base models encode
grammaticality, so both endpoints of the $\lambda$ axis are fluent language models
and the segment between them lies inside the manifold of well-formed text.

\paragraph{A $\psi$ threshold was derived and does not hold.} The same model was
used to derive a crossover $\psi^{*} \approx 0.71$--$0.83$ below which reference
subtraction should lose to temperature, and three earlier drafts compared measured
$\psi$ against it. $\psi^{*}$ is not a constant: holding the model's other
parameters fixed and varying only the matched-diversity target moves it from 0.000
to 1.000, and the published value is the most demanding point on that curve. The
model is also off-regime --- its aligned distribution collapses 47-fold against its
reference, where the measured collapse is 1.5--2.2$\times$ under continuation
scoring (1.51--3.22$\times$ under prompt-final) --- and rebuilt at measured
peakedness and collapse it yields $\psi^{*}$ between 0.000 and 0.424, below several
measured values. \textbf{We withdraw the threshold and make no decision-theoretic
claim.} $\psi$ is reported as an absolute magnitude throughout.

The assumption behind both failures --- quality independent of typicality --- is
common in this literature and reverses conclusions.

\section{Related work}

\paragraph{Decoding-time model composition} is mature and has been used almost
exclusively to add capability. Contrastive decoding \cite{li2023} maximises
expert-minus-amateur log-likelihood to suppress generic output. Proxy-tuning
\cite{liu2024proxy} adds a small-model tuning delta to a large base model. DeRa
\cite{liu2024dera} proves the geometric-mixture structure and sweeps
$\lambda \in [0, 2]$. None evaluates against diversity metrics and none explores
$\lambda < 0$.

\paragraph{The typicality-bias literature} proposes prompt-level interventions:
verbalized sampling, diverse personas, chain-of-thought differentiation
\cite{zhang2025}. These require the user to want diversity, which matters given that
the collective-diversity loss Doshi and Hauser \cite{doshi2024} identify is a social
dilemma no individual has an incentive to solve. Related work on preference collapse
\cite{xiao2024} shows minority preferences being virtually disregarded under
standard alignment objectives.

\paragraph{Stage-resolved analysis.} Kirk et al.\ \cite{kirk2024} compare SFT against
RLHF on generalisation and diversity but do not decompose a single pipeline's
entropy budget across its stages. We have not found a cross-pipeline comparison of
such a decomposition, nor a stage comparison that reports the position convention of
Section~\ref{sec:conv}.

\section{Limitations}
\label{sec:limits}

\begin{enumerate}
\item \textbf{Four pipelines, two base families.} The shared-shift result replicates
within two lineages that use related instruction mixtures. The OLMo 7B and 13B
share one release, so their agreement is a within-release scale comparison, not an
independent replication. The unit of inference is the prompt: 36 prompts in six
task families, however many positions each contributes.
\item \textbf{The lineage test rules out ancestry, not data.} Both families train on
T\"ulu-3-style mixtures, so shared training data remains a live explanation.
\item \textbf{$\beta$ was measured at the 1B, 7B and T\"ulu-8B}, on 400 sampled
positions, but not at the 13B. The 1B preference stage is the one place it is
$\approx 0$. The $R^2$ and Shapley results in Section~5.1 cover all four models.
\item \textbf{``Typicality'' is operationalised} as a linear and piecewise-linear
function of reference log-probability. It is one operationalisation of a claim the
source literature does not state formally. Section~5.2 shows that two analyst
choices --- linear versus non-linear, and weighted versus unweighted --- each
reverse a conclusion.
\item \textbf{One judge.} GPT-2 perplexity measures fluency, not whether an answer is
good, and it rewards predictable text, biasing against diverse conditions. Human
evaluation is absent and is the limitation we would most want addressed.
\item \textbf{The $\gamma$ sweep is underpowered} at five prompts and resolves
nothing.
\item \textbf{The 1B $\lambda$ null is unexplained.}
\item \textbf{Scale.} 135M, 1B, 7B, 8B, 13B. Nothing here speaks to frontier scale.
\end{enumerate}

\section{Conclusion}

Across four pipelines, alignment's diversity cost is carried by a small part of the
alignment shift. A typicality-aligned component accounting for at most about a
third of a stage's shift carries nearly all of its entropy loss under the
probability-weighted fit, at both stages, in every model measured and across the
specifications we tried; under an unweighted fit its share of the loss is smaller
and, at the 1B's instruction tuning, unstable in sign. Larger
than that component, though itself a minority of the shift, is a shared direction:
measurable at every stage of every model, surviving a partial correlation against
reference typicality, with null controls at 0.008 and 0.004. Its context-invariance
depends on stage and scale together, and in opposite directions: within OLMo,
instruction tuning rises with scale (0.092, 0.320, 0.445 at 1B, 7B and 13B) while
the preference stage falls (0.558, 0.489, 0.386). It is moderately
context-invariant everywhere --- $\rreg$ between 0.26 and 0.56 --- except for
instruction tuning at the 1B. Squared, that is at most roughly 0.31 of the shift ---
the square of a mean correlation, so an approximation rather than a bound --- and it is
not a context-independent direction so much as a partially shared one. Most of each
stage's shift is accounted for by neither axis, we do not characterise what the
remainder is, and it is not inert --- at instruction tuning in the 1B it raises
entropy by enough to offset a third of the loss (Section~5.1).

Those findings suggest a hypothesis about larger scales, which this work does not
test: an operator on the reference axis acts on the axis that carries the entropy,
and at 7--8B the entropy-relevant part of the preference shift is already close to a
temperature, so there would be little for such an operator to recover that
temperature does not. It is a hypothesis and not an account of what we measured. No
decoding-time intervention was run above 1B. It cannot be the explanation of the 1B
null, where the preference stage's $\beta \approx 0$ and $\lambda$ acts mostly on
instruction tuning (Section~6.3); that null stays unexplained, and the same operator
beat temperature threefold at 135M. The results also raise a
question this work does not answer. The field measures alignment's effect through
diversity metrics, and most of what alignment does to the next-token distribution
does not appear in them.

Underneath sits a methodological result we did not expect to be consequential. The
same T\"ulu-3-8B checkpoints yield preference-stage shares of 29\% and 53\%
depending on where the distribution is scored, and the 1B model's conventional
measurement sits 18--28 points from every alternative. Two of this work's own
earlier conclusions were reversed by that choice alone. Stage attribution is not
reproducible without stating it.

\appendix

\section{Reproduction}

\begin{center}
% p-columns: both sides grew long once the release was completed, and an ll table
% overflows the text block as soon as either does.
\begin{tabular}{p{0.30\textwidth} p{0.60\textwidth}}
\toprule
file & contents \\
\midrule
\texttt{dump\_logprobs\_mps.py} & per-checkpoint log-prob dumps \\
\texttt{analyse\_psi\_family.py} & stage-resolved $\psi$ and entropy, bootstrap over prompts \\
\texttt{register\_v2.py} & shared-shift geometry, both arms, null controls \\
\texttt{decomp.py} & counterfactual Shapley entropy decomposition \\
\texttt{decomp\_resid.py} & non-linear typicality, cross-fitted shared directions \\
\texttt{robustness.py} & specification grid for Sec.~\ref{sec:minority} (weighting $\times$ plausible threshold $\times$ held-out shared direction), R$^2$ and Shapley shares under the same fit, per-family and leave-one-family-out shares (exploratory, v1.0) \\
\texttt{check\_tokenizers.py} & token-ID-level tokeniser identity verification \\
\texttt{grind.py}, \texttt{matched6.py} & generation sweep and matched-diversity scoring \\
\texttt{gen\_scale.py} (mode \texttt{ctrl}) & shuffled-delta direction control \\
\texttt{exp1e\_peakedness.py} & the probe that withdrew the $\psi^{*}$ threshold (Sec.~\ref{sec:failures}) \\
\texttt{pilot135m\_port/} & 135M $\gamma$ sweep and $\beta$ measurement, ported \\
\texttt{pilot-135m/} & the original 135M pilot, including the withdrawn threshold and the source of the 3.00$\times$ figure \\
\texttt{logs/} & run scripts and evaluation outputs behind the reported numbers \\
\texttt{logs/inline\_checks.py} & reproduces the figures originally computed inline rather than by a script \\
\bottomrule
\end{tabular}
\end{center}

Six pre-registrations, each hashed before its data existed, are released with the code
at \url{https://github.com/amitd01/axisalignment-paper}: the 13B scale test, the register test, the lineage test, the entropy
decomposition, the residual test and the 135M $\gamma$ sweep. Two carry dated notes
appended after the fact; the hashes cover the original text and the notes state what
changed and when.

\begin{center}
\scriptsize
\begin{tabular}{p{0.13\textwidth} p{0.09\textwidth} p{0.14\textwidth} p{0.19\textwidth} p{0.19\textwidth} p{0.13\textwidth}}
\toprule
pre-registration & written (IST) & already seen when written & note appended & outcome & standing \\
\midrule
13B scale test & 2026-10-13 20:20 & 1B and 7B prompt-final shares & 21:21, before any 13B dump: comparison basis moved to continuation positions & trend continues, +15.6 [+12.7, +18.6] (Sec.~\ref{sec:stages}) & confirmatory, amended \\
register test & 2026-10-13 20:22 & no register data & --- & $\rho_{\mathrm{reg}} > \psi$, all four (Sec.~\ref{sec:exceeds}) & confirmatory \\
lineage test & 2026-10-13 22:20 & same-lineage value $-$0.19 & --- & $-$0.194, transfer (Sec.~\ref{sec:lineage}) & confirmatory \\
entropy decomposition & 2026-10-13 22:25 & stage shares; no decomposition & --- & head $>$ 50\%: not supported; typicality $>$ shared: holds (Sec.~\ref{sec:cost}) & confirmatory \\
residual test & 2026-10-13 23:13 & decomposition shares & before the run: position-specific claim withdrawn & non-linear claim too strong at 1B, T\"ulu (Sec.~\ref{sec:fitdep}) & confirmatory for the non-linear question; $\varepsilon$ share descriptive \\
135M $\gamma$ sweep & 2026-10-13 23:28 & published 135M scores & --- & unresolved (Sec.~\ref{sec:refsub}) & confirmatory, underpowered \\
\bottomrule
\end{tabular}
\end{center}

Not pre-registered, and reported as exploratory: the variance shares (R$^2$) of
Sec.~\ref{sec:minority} and the specification grid in \texttt{robustness.py}. The adoption of
base-written continuations as the primary convention for stage shares (Sec.~\ref{sec:fourconv}) is
the 13B row's appended note.

Seeds: 20261012 (1B), 20260803 (135M). $N = 6$ samples, 32 new tokens, top-$p$ 0.95
held constant. Fixed seeds do not guarantee bit-identical output across hardware
\cite{tml2025}; batch-invariant kernels were not attempted.

\paragraph{Consistency check.} $T = 1.0$, $\lambda = +1.0$ and $\gamma = 0$ are the
same distribution by construction and returned identical metrics on the original
cloud run (1B: Jaccard 0.0746, ppl 58.8). It was not re-run on the hardware used for
Sections 3--5.

\paragraph{Metrics.} Entropy is the diversity measure throughout Sections 2--5.
Section~\ref{sec:decode} uses mean pairwise Jaccard; distinct-$n$ saturates at 1.000
by $T = 1.6$ at these scales and is reported only for comparability with prior work.
Perplexity is a median across samples.

\begin{thebibliography}{99}

\bibitem{li2023}
X.~L. Li et al.
\newblock Contrastive Decoding: Open-ended Text Generation as Optimization.
\newblock arXiv:2210.15097, 2023.

\bibitem{liu2024proxy}
A.~Liu et al.
\newblock Tuning Language Models by Proxy.
\newblock arXiv:2401.08565, COLM 2024.

\bibitem{liu2024dera}
T.~Liu et al.
\newblock Decoding-time Realignment of Language Models.
\newblock arXiv:2402.02992, ICML 2024.

\bibitem{kirk2024}
R.~Kirk et al.
\newblock Understanding the Effects of RLHF on LLM Generalisation and Diversity.
\newblock arXiv:2310.06452, ICLR 2024.

\bibitem{zhang2025}
S.~Zhang et al.
\newblock Verbalized Sampling: How to Mitigate Mode Collapse and Unlock LLM
Diversity.
\newblock arXiv:2510.01171, 2025.

\bibitem{xiao2024}
J.~Xiao et al.
\newblock On the Algorithmic Bias of Aligning Large Language Models with RLHF:
Preference Collapse and Matching Regularization.
\newblock arXiv:2405.16455, 2024.

\bibitem{doshi2024}
A.~R. Doshi and O.~P. Hauser.
\newblock Generative AI enhances individual creativity but reduces the collective
diversity of novel content.
\newblock \emph{Science Advances}, 10(28), 2024.

\bibitem{sourati2026}
Z.~Sourati, A.~H. Ziabari and M.~Dehghani.
\newblock Flattening the cognitive landscapes: how generative AI reshapes human
thought.
\newblock \emph{Trends in Cognitive Sciences}, 2026.

\bibitem{olmo2}
Allen Institute for AI.
\newblock OLMo 2: base, SFT, DPO and Instruct checkpoints.
\newblock \url{https://huggingface.co/allenai}, 2024--2025.

\bibitem{tulu3}
N.~Lambert et al.
\newblock T\"ulu 3: Pushing Frontiers in Open Language Model Post-Training.
\newblock arXiv:2411.15124, 2024.

\bibitem{tml2025}
Thinking Machines Lab.
\newblock Defeating Nondeterminism in LLM Inference.
\newblock 2025.

\end{thebibliography}

\end{document}
